\chapter{From glTF File to Model: the Import Core}
\label{ch:gltf-import-core}

CNA does not have a small ``glTF reader'' hidden inside \cnaclass{ContentManager}. It has a
shared import library, \texttt{CNA::Internal::GltfImport}, used by two independent front ends:
the runtime \texttt{.gltf}/\texttt{.glb} path and the offline
\texttt{cna\_tool\_gltf\_to\_cnj} converter. This distinction matters. Parsing and semantic
extraction are shared; turning the extracted records into live graphics objects or persistent
CNJ sidecars is not.

The shared core is about 2,500 lines across \texttt{GltfImportCore.hpp} and
\texttt{GltfImportCore.cpp}. It delegates container parsing and accessor decoding to the
vendored cgltf~1.15 library. Exactly one translation unit defines
\texttt{CGLTF\_IMPLEMENTATION}; the stb image implementations in the same translation unit
are declared static so they cannot collide with another copy linked through SDL image support.
That is a build invariant, not incidental preprocessor decoration.

\section{The parser boundary}

Both JSON \texttt{.gltf} and binary \texttt{.glb} enter through the same
\texttt{cgltf\_parse\_file} path and must declare \texttt{asset.version} exactly
\texttt{"2.0"}. Buffers are then loaded relative to the source path. Embedded buffer views,
base64 data URIs, and external percent-decoded file URIs are all supported. Images use the same
three storage forms and are exposed as memory blocks; the runtime front end decodes those blocks
without staging temporary files.

This is not equivalent to full glTF validation. CNA never calls \texttt{cgltf\_validate}, and
does not inspect \texttt{extensionsRequired} or \texttt{extensionsUsed}. A syntactically
parseable file that declares an unsupported required extension can therefore continue until a
later semantic limitation either rejects it or silently loses information. The importer's own
named checks remain valuable---accessor shape, buffer span, index range, primitive mode, and
version failures generally identify the offending value---but they are not a substitute for a
validator in an asset-production pipeline.

\paragraph{Practical rule.}
Run the Khronos validator before shipping third-party content. A successful CNA load proves
that the code paths it used were acceptable; it does not prove that the complete document is
spec-conformant or that every required extension was honoured.

\section{Scene graph first, meshes second}

\texttt{BuildSceneGraph} selects the document's default scene, falling back to the first scene,
then walks it iteratively in parent-before-child order. Index~0 is a synthetic identity node
named \texttt{Root}; each reachable glTF node follows exactly once. The core asks cgltf for the
local matrix, converts it to CNA's row-vector convention, and composes

\[
  W_i = L_i W_{parent(i)}.
\]

The conversion copies the affine basis into CNA's matrix convention. It does not apply an axis
or handedness conversion, and that is correct: both glTF and the XNA conventions used here are
right-handed, use $+Y$ as up, $-Z$ as forward, and place the UV origin at the top left. The
only handedness value imported separately is a tangent's fourth component.

Each flattened scene node becomes a \cnaclass{ModelBone}, index-for-index after the synthetic
root. Vertex positions remain in mesh-local space. Consequently one glTF mesh instanced by two
nodes yields two \cnaclass{ModelMesh} placements with distinct parent bones instead of one
destructively pre-transformed vertex buffer. For rigid geometry the mesh is parented to its
instancing node. For skinned geometry the node still exists, but the mesh is parented to the
synthetic root because skin matrices already account for the mesh-node coordinate space; doing
both would apply the transform twice.

Nodes unreachable from the selected scene do not enter this graph. If no scene node references
any mesh, the importer has a compatibility fallback that exposes every mesh at the identity
root. That fallback is useful for incomplete authoring exports, but it is not the same semantic
result as a fully authored scene.

\section{Groups are an ownership boundary}

\texttt{CollectMeshGroups} partitions scene instances into one group for each distinct
\texttt{skin} pointer plus one group for unskinned content. A group record contains the source
node, mesh, flattened scene-node index, composed world transform, and whether the instance is
skinned. This division is the unit that later becomes one model output.

The offline converter writes every group. A file containing static scenery and two characters
can therefore produce three CNJ models, with stable suffixes for the static group and named
skins. The runtime \texttt{Load<Model>} path deliberately returns only
\texttt{groups.front()}. It does not return a scene object and does not warn that later groups
were discarded. Multi-character or mixed static/skinned files should be converted offline or
split at authoring time.

The two front ends also differ in scale control. The core's \texttt{unitScale} multiplies
translations---including vertices, bind translations, animation translations and tangents,
morph position deltas, inverse binds, and ancestor terms---but never rotations or dimensionless
scale factors. The runtime path fixes it at~1.0; only the command-line converter exposes a
caller-supplied value.

\section{What extraction preserves}

For each triangle primitive, the core produces an explicit semantic record before any graphics
object is allocated. It decodes positions, normals, tangents, texture coordinates, colours,
joint indices and weights; validates every index; synthesises a sequential index stream for a
non-indexed primitive; and computes tangents when required by a normal-mapped material. Sparse
attributes and sparse indices, including an accessor with no base buffer view, are handled.

Materials are less complete. PBR effects are selected only when a normal or
metallic--roughness \emph{map} is present. Metallic, roughness, and emissive factors survive
only after that selection; \texttt{baseColorFactor}, alpha mode/cutoff, double-sidedness,
sampler state, and colour-space intent are not represented. A factor-only metallic--roughness
material therefore becomes a white \cnaclass{BasicEffect}. This is the open D7 defect, not a
supported approximation.

Exactly four extension families have production handling:

\begin{itemize}
  \item \texttt{KHR\_texture\_transform}, for the base-colour texture, baked into the one UV
        stream;
  \item \texttt{KHR\_materials\_emissive\_strength}, on the PBR path;
  \item \texttt{KHR\_lights\_punctual}, reduced to at most three directional-light
        approximations; and
  \item \texttt{KHR\_draco\_mesh\_compression}, only when CNA was compiled with libdraco.
\end{itemize}

Transmission, variants, cameras, glTF sampler choices, and negative-determinant winding repair
are absent. A PBR map that selects a second UV set is detected through
\texttt{pbrUv2Mismatch}; the CLI warns, while the runtime currently computes and then hides the
flag. Keep ``parsed'', ``represented'', and ``drawn'' as three separate claims.

\section{Failure and lifetime boundaries}

The core throws \texttt{std::runtime\_error} for several semantic rejections, including an
unsupported primitive topology. The runtime reader does not normalize all of these into
\cnaclass{ContentLoadException}; a caller catching only the latter can miss a glTF import
failure even though it derives from \texttt{std::runtime\_error}. This exception asymmetry is
part of the current contract.

On success, the runtime front end creates model bones, meshes, parts, buffers, effects,
textures, skinning records, and morph records in one private ownership bundle, then attaches
that bundle to the otherwise raw-pointer-heavy \cnaclass{Model}. Images are cached by source
image pointer for that one import. The result is self-retaining, but repeated primitives still
receive their own mesh-part graphics resources unless the front end explicitly shares them.

The next four chapters follow the same pipeline at progressively sharper boundaries: exact GPU
bytes in Chapter~\ref{ch:gpu-packing}, skeletal coordinate spaces in
Chapter~\ref{ch:skinning-animation}, persistent CNJ output in
Chapter~\ref{ch:cnj-model-toolchain}, and the conformance evidence in
Chapter~\ref{ch:gltf-conformance}.
